\chapter{Implementation and Development}

\section{Technology Stack Implementation}

\subsection{Frontend Implementation}
The frontend was developed using React 19, providing a modern Single Page Application experience:

\begin{itemize}
    \item \textbf{React 19}: Latest version with improved performance and concurrent rendering
    \item \textbf{React Router DOM 7}: Client-side routing for seamless navigation
    \item \textbf{Supabase JS Client}: Authentication handling on the frontend
    \item \textbf{Custom CSS}: Aesthetic styling for the virtual study environment
\end{itemize}

Key frontend components implemented:
\begin{itemize}
    \item \texttt{App.js} - Root component with routing configuration
    \item \texttt{Dashboard.jsx} - Learning progress tracker with gamification
    \item \texttt{StudyRoom.jsx} - Virtual study environment with hotspots
    \item \texttt{Calendar.jsx} - Monthly calendar with event management
    \item \texttt{Notes.jsx} - Notes management with image support
    \item \texttt{Journal.jsx} - Daily journal with mood tracking
    \item \texttt{PomodoroTimer.jsx} - Focus timer with break management
    \item \texttt{QuizModal.jsx} - Quiz interface for self-assessment
\end{itemize}

\subsection{Backend Implementation}
The backend was built using Node.js with Express.js framework:

\begin{itemize}
    \item \textbf{Express.js 4.18}: RESTful API framework
    \item \textbf{pg (node-postgres)}: PostgreSQL database client
    \item \textbf{Supabase}: Authentication service integration
    \item \textbf{pdfjs-dist}: PDF parsing for quiz generation
    \item \textbf{Helmet}: Security middleware
    \item \textbf{Morgan}: Request logging
    \item \textbf{CORS}: Cross-origin resource sharing
\end{itemize}

Backend route modules implemented:
\begin{itemize}
    \item \texttt{auth.js} - Authentication endpoints
    \item \texttt{tasks.js} - Task CRUD operations
    \item \texttt{notes.js} - Notes management
    \item \texttt{journal.js} - Journal entries
    \item \texttt{calendar.js} - Calendar events
    \item \texttt{files.js} - File upload/management
    \item \texttt{focus-sessions.js} - Pomodoro session tracking
    \item \texttt{progress.js} - Learning progress and XP
    \item \texttt{quiz.js} - Quiz generation from documents
    \item \texttt{trash.js} - Soft delete management
\end{itemize}

\subsection{Database Implementation}
PostgreSQL database with 12 tables:

\begin{table}[H]
\centering
\begin{tabular}{|l|p{8cm}|}
\hline
\textbf{Table} & \textbf{Purpose} \\
\hline
users & User profiles, XP, streaks, last activity \\
tasks & Task items with priority, status, progress \\
notes & Quick notes with optional images \\
journal\_entries & Daily journal with mood tracking \\
calendar\_events & Events with multi-day support \\
files & Uploaded documents for quiz generation \\
focus\_sessions & Pomodoro session records \\
learning\_sections & 3 learning phases \\
learning\_levels & 15 levels (5 per phase) \\
progress & User progress tracking \\
trash & Soft-deleted items for recovery \\
schedules & Time slot allocations \\
\hline
\end{tabular}
\caption{Database Tables Implemented}
\end{table}

\section{Core Features Implementation}

\subsection{Dashboard and Learning Progress}
The dashboard implements a gamification system with:
\begin{itemize}
    \item \textbf{3 Phases}: Foundations \& Basics, Building Momentum, Peak Mastery
    \item \textbf{15 Levels}: 5 levels per phase with progressive unlocking
    \item \textbf{XP System}: Points earned from completing quizzes and tasks
    \item \textbf{Streak Tracking}: Consecutive days of activity
    \item \textbf{Productivity Score}: Calculated from completed activities
\end{itemize}

\subsection{Virtual Study Room}
The Study Room provides an immersive environment with:
\begin{itemize}
    \item Video background for ambient atmosphere
    \item Interactive hotspots for:
    \begin{itemize}
        \item Calendar access
        \item Todo list
        \item Bookshelf (file management)
        \item Spotify integration
        \item Trash management
    \end{itemize}
    \item Pomodoro timer integration
    \item Focus mode with distraction-minimized view
\end{itemize}

\subsection{Pomodoro Timer Implementation}
The focus timer includes:
\begin{itemize}
    \item Configurable session duration (default 25 minutes)
    \item Start, pause, and resume functionality
    \item Automatic break reminders (5 minutes)
    \item Task association for progress tracking
    \item Session history and analytics
\end{itemize}

Session lifecycle states:
\begin{enumerate}
    \item \texttt{START} - Session initialized
    \item \texttt{IN\_PROGRESS} - Timer running
    \item \texttt{PAUSED} - Timer temporarily stopped
    \item \texttt{COMPLETED} - Session finished successfully
    \item \texttt{CANCELLED} - Session ended early
\end{enumerate}

\subsection{Quiz Generation System}
The quiz module implements automatic question generation:

\textbf{Algorithm Steps:}
\begin{enumerate}
    \item Extract text from uploaded PDF/TXT file using pdf.js
    \item Split text into sentences (minimum 45 characters)
    \item Build word frequency map excluding stop words
    \item Identify key terms from each sentence
    \item Create fill-in-the-blank questions
    \item Generate distractors from top frequent words
    \item Return 5 MCQs with shuffled options
\end{enumerate}

\subsection{Calendar System}
The calendar implements:
\begin{itemize}
    \item Monthly grid view
    \item Event types: Reminder, Exam, Deadline, Study Session, Other
    \item Color coding by event type
    \item Multi-day event spanning
    \item Drag-and-drop rescheduling
    \item Task synchronization (tasks appear on target dates)
    \item Side panel for quick access
\end{itemize}

\subsection{Notes Module}
Notes implementation features:
\begin{itemize}
    \item Rich text content storage
    \item Image support (up to 5 images per note)
    \item Search functionality
    \item Random pastel color coding
    \item Soft delete with trash recovery
\end{itemize}

\subsection{Journal Module}
Daily journal features:
\begin{itemize}
    \item One entry per day constraint
    \item 8 mood options (Happy, Excited, Sad, Tired, Thinking, Cool, Party, Peaceful)
    \item Automatic word count
    \item Calendar-based entry navigation
\end{itemize}

\subsection{Focus Timeline Analytics}
Analytics views implemented:
\begin{itemize}
    \item \textbf{Daily View}: Hourly breakdown of focus sessions
    \item \textbf{Weekly View}: Week-by-week summary
    \item \textbf{Monthly View}: Calendar heat map visualization
    \item Statistics: Total focus time, sessions completed, productivity metrics
\end{itemize}

\subsection{Export Functionality}
Report export capabilities:
\begin{itemize}
    \item \textbf{CSV Export}: Tabular data for spreadsheet analysis
    \item \textbf{PDF Export}: Formatted reports for documentation
    \item Data includes: Tasks, focus sessions, progress metrics
\end{itemize}

\section{Data Synchronization}
The system maintains real-time synchronization using custom events:

\begin{itemize}
    \item \texttt{studybug:tasks-updated} - Task changes
    \item \texttt{studybug:calendar-updated} - Calendar modifications
    \item \texttt{studybug:focus-completed} - Session completion
\end{itemize}

Components listen for these events and refresh their data automatically, ensuring consistency across the Calendar, Todo, Pomodoro, and Dashboard views.

\section{Security Implementation}
\begin{itemize}
    \item JWT-based authentication with Supabase
    \item User data isolation (users can only access their own data)
    \item Parameterized SQL queries to prevent injection
    \item Helmet middleware for security headers
    \item Input validation on all API endpoints
    \item Session expiry validation
\end{itemize}

\section{Development Statistics}

\begin{table}[H]
\centering
\begin{tabular}{|l|c|}
\hline
\textbf{Metric} & \textbf{Value} \\
\hline
Total Files Created & 100+ \\
Frontend Components & 20+ \\
API Endpoints & 35+ \\
Database Tables & 12 \\
CSS Files & 15+ \\
Total Commits & 52 \\
Pull Requests Merged & 13 \\
\hline
\end{tabular}
\caption{Development Statistics}
\end{table}
